\chapter{脑电采集硬件系统设计与实现}
\label{cha:hardware}

为满足低通道脑机接口系统对便携化、低功耗和无线化采集平台的需求，本文设计并实现了一套面向双通道脑电采集的嵌入式硬件系统。该系统以脑电前端采集链路为核心，结合低功耗主控芯片、蓝牙无线传输链路以及配套上位机软件，形成了从信号采集、参数配置、数据传输到实时可视化显示的完整闭环。与仅面向离线算法研究的数据集环境不同，本章关注的是面向真实采集场景的系统实现问题，即如何在有限板级资源与无线带宽约束下完成双通道脑电信号的稳定获取，并为后续算法验证与系统实验提供可复用的硬件平台。

本章将从总体架构、关键电路、固件流程以及系统联调验证四个层面展开论述。其中，总体架构部分重点说明系统模块划分及其协同关系；电路设计部分围绕脑电采集前端、主控与无线传输链路以及供电接口设计进行分析；固件部分说明采样、滤波、缓冲与分帧传输的实现逻辑；最后结合当前已完成的联调结果，对系统的闭环可用性与实现特征进行分析。

\section{硬件系统总体架构设计}

\subsection{模块功能分工}

本文所设计的硬件系统由脑电采集前端模块、主控与无线传输模块、供电与接口模块以及配套上位机软件组成，其核心目标是在保证系统便携性和可配置性的前提下，实现双通道脑电信号的连续采集与实时显示。系统的主数据链路为：电极采集到的模拟脑电信号首先进入前端采集芯片完成模拟调理与增益控制，随后由主控芯片进行模数转换、缓存管理与数字滤波，并通过 BLE 链路发送至上位机，由上位机完成帧解析、波形显示、频带能量估计和参数配置。除脑电主链路外，系统还预留了基于 I\textsuperscript{2}C 总线的基础生理信号扩展接口，用于实现心率与血氧等辅助生理信息的采集与显示。

从系统级角度看，该平台并非单一的采样板卡，而是一个围绕双通道脑电采集构建的完整实验终端。一方面，前端采集芯片负责弱脑电信号的模拟耦合与放大，为后续数字化提供稳定输入；另一方面，低功耗主控芯片同时承担参数控制、数据打包和无线通信职能，使系统在板级面积、功耗和复杂度之间取得平衡。各模块分工如表\ref{tab:c3_module_split}所示。

\begin{table}[htbp]
    \centering
    \caption{双通道脑电采集硬件系统的模块划分与主要职责}
    \label{tab:c3_module_split}
    \begin{tabular}{p{3.2cm}p{3.4cm}p{7.2cm}}
        \hline
        模块 & 主要器件/载体 & 主要职责 \\
        \hline
        脑电采集前端模块 & KS1092、输入接口与保护网络 & 完成双通道脑电模拟信号接入、前端增益配置、导联状态检测，并向主控输出可采样的模拟信号 \\
        主控与无线传输模块 & nRF52832、片上 SAADC、BLE 链路 & 完成双通道采样、参数配置、数字滤波、缓冲区管理和无线数据发送 \\
        供电与接口模块 & Type-C、充电管理与稳压电路、电池接口 & 提供便携式供电能力，支撑板级上电、自锁和充电管理 \\
        配套上位机软件 & Qt6 上位机程序 & 完成设备连接、帧解析、实时波形显示、频带能量估计与参数下发 \\
        生理信号扩展模块 & I\textsuperscript{2}C 扩展链路 & 提供心率/血氧等基础生理参数的扩展采集与显示能力 \\
        \hline
    \end{tabular}
\end{table}

从板级结构看，该系统采用单板集成式实现方案，将脑电前端、主控、充电与接口等功能集中在一块紧凑 PCB 上，以降低系统体积并减少跨板连接带来的不稳定因素。图\ref{fig:c3_pcb_3d}给出了硬件平台的三维结构预览。可以看出，系统顶部集中布置了 Type-C 接口、电池连接器、主控与外围晶振、按键以及脑电前端芯片等核心器件，整体结构紧凑，便于佩戴式或近体式部署。该布局使脑电采集、无线传输和电源管理三个关键功能在同一板级平台上形成闭环，为后续的移动化实验与算法联调提供了硬件基础。

\begin{figure}[htbp]
    \centering
    \begin{minipage}[b]{0.48\linewidth}
        \centering
        \includegraphics[width=\linewidth]{C3-3_PCB_3D_Top.png}
    \end{minipage}
    \hfill
    \begin{minipage}[b]{0.48\linewidth}
        \centering
        \includegraphics[width=\linewidth]{C3-4_PCB_3D_Bottom.png}
    \end{minipage}
    \caption{双通道脑电采集硬件平台的 PCB 三维结构预览}
    \label{fig:c3_pcb_3d}
\end{figure}

\subsection{接口协议设计}

为适应低功耗无线采集场景，本文采用 Nordic UART Service（NUS）作为设备与上位机之间的通信承载协议。该方案的优点在于实现复杂度较低、兼容性较好，并可同时支持连续数据上传与文本命令交互。考虑到系统既需要持续发送脑电数据帧，又需要支持参数配置、状态查询与错误回执，协议设计采用了“二进制连续数据帧 + ASCII 文本命令/回执”的混合方式。其中，脑电波形数据采用紧凑的二进制帧结构发送，以提高有效负载利用率；命令控制与状态回执则使用 ASCII 文本消息，便于上位机解析与调试。

在脑电数据链路中，每帧数据以固定魔数 `0xE0` 作为起始标识，用于与文本回执消息区分；帧头中进一步包含序列号与样本数量字段，随后按顺序附带多个双通道采样点。每个采样点包含两个通道的原始 ADC 值以及设备侧时间戳。该结构既能满足上位机对连续波形的实时显示需求，又保留了基础时序信息，便于后续进行丢包检测与链路分析。其帧结构如表\ref{tab:c3_ble_frame}所示。

\begin{table}[htbp]
    \centering
    \caption{脑电数据的 BLE 二进制帧结构}
    \label{tab:c3_ble_frame}
    \begin{tabular}{p{3.0cm}p{2.0cm}p{2.5cm}p{6.0cm}}
        \hline
        字段 & 长度/字节 & 数据类型 & 说明 \\
        \hline
        Magic & 1 & `uint8` & 固定为 `0xE0`，用于区分脑电二进制帧与文本消息 \\
        Sequence ID & 2 & `uint16` & 数据包序号，用于检测丢包与乱序 \\
        Sample Count & 1 & `uint8` & 当前数据帧包含的采样点数量 \\
        Channel 0 ADC & 2 & `int16` & 通道 0 原始采样值 \\
        Channel 1 ADC & 2 & `int16` & 通道 1 原始采样值 \\
        Timestamp & 4 & `uint32` & 设备侧时间戳，用于时序对齐与状态分析 \\
        \hline
    \end{tabular}
\end{table}

除脑电二进制帧外，系统还支持基于文本命令的参数配置与状态查询，包括后端增益设置、滤波模式设置、缓存状态查询、导联脱落状态上报以及脑电前端复位等功能。与纯被动显示的数据终端不同，该协议设计使上位机能够参与采集链路的闭环控制，即在保持采样主链路轻量化的同时，为实验过程中的调参与故障恢复保留必要接口。这一设计对于后续的多轮联调和不同实验条件下的快速配置具有重要意义。

\subsection{配套上位机软件闭环设计}

在本系统中，上位机并非独立于硬件之外的附属调试工具，而是采集闭环中的重要组成部分。其主要职责包括完成 BLE 设备发现与连接、订阅数据通知、区分二进制脑电帧与文本回执消息、实时显示双通道波形、估计频带能量分布以及向硬件下发增益和滤波配置命令。通过将这些功能整合到统一界面中，系统能够在采集过程中同时完成状态观测与参数控制，从而提高联调效率。

图\ref{fig:c3_ui}给出了配套上位机界面示意。界面左侧为设备连接与状态区，可显示设备连接状态与导联状态；中部为双通道原始脑电实时波形；下方进一步显示频带能量分布和部分生理参数信息；右侧则提供后端增益调节与滤波模式设置功能。该界面布局的设计思路并不是追求复杂交互，而是围绕“采集状态是否正常、波形是否连续、参数是否可配置”这三个核心问题服务于硬件系统本身。因此，在第 3 章中，上位机主要被视为硬件平台的配套实现部分，而其在实验中的实际表现与效果分析将放在后续章节中进一步展开。

\begin{figure}[htbp]
    \centering
    \includegraphics[width=0.95\linewidth]{C3-5_Desktop_UI_Screenshot.png}
    \caption{配套上位机软件界面}
    \label{fig:c3_ui}
\end{figure}

\section{核心硬件电路设计}

系统的硬件电路设计围绕“弱脑电信号稳定获取、低功耗主控处理以及无线链路可靠传输”三项要求展开。考虑到脑电信号幅值较低且易受环境干扰影响，前端电路需要兼顾输入保护、导联连接稳定性与可调增益能力；而主控部分不仅要承担模数转换与通信任务，还需要在有限资源条件下支撑蓝牙协议栈和实时采样任务。基于上述考虑，本文采用 KS1092 作为脑电前端采集芯片，采用 nRF52832 作为主控与无线传输核心，并配合电源管理与外围接口电路形成完整硬件平台。

\subsection{脑电采集前端电路设计}

脑电采集前端电路是整套系统中最关键的模拟链路，其作用在于将来自电极接口的微弱脑电信号稳定耦合至后级采样链路。图\ref{fig:c3_front_end}给出了脑电前端局部原理图。从电路结构可以看出，系统通过多针接口引出双通道输入与参考/驱动信号，前端输入侧配置了保护网络，以增强实际采集条件下对静电与外部瞬态干扰的承受能力。双通道脑电信号分别进入 KS1092 的差分输入端，形成两条并行采集链路。

在功能实现上，KS1092 负责完成前端增益配置、导联脱落检测以及与主控之间的配置交互。芯片的配置接口通过 `SCLK`、`SDI`、`SDO` 与 `CS` 与主控侧相连，`RESET` 与 `EN` 引脚用于上电初始化与工作状态控制，`LDF` 信号则用于向主控反馈电极接触状态。该设计使脑电前端不再是固定参数的黑盒模块，而是一个可被主控动态配置并具备基础自检能力的采集单元。结合固件实现，系统支持多个离散增益档位配置，从而可针对不同被试、不同电极接触条件和不同实验任务进行前端调节。

需要指出的是，当前系统采用的是“模拟前端 + 主控侧 ADC 采样”的实现路径，即 KS1092 完成前端模拟调理后，将输出信号送入主控芯片的模数转换通道，而不是直接将数字化脑电数据从前端芯片侧输出。这一路径在硬件资源和板级复杂度上更为可控，同时也为系统保留了在主控侧执行滤波、缓存与打包策略的灵活性。对于本研究所关注的低通道便携式系统而言，这种设计兼顾了实现复杂度与实验可扩展性。

\begin{figure}[htbp]
    \centering
    \includegraphics[width=0.92\linewidth]{C3-1_Front_End_Schematic.png}
    \caption{脑电采集前端局部原理图}
    \label{fig:c3_front_end}
\end{figure}

从原理图中还可以看到，前端电路中引入了与参考驱动相关的反馈网络。该网络用于改善输入参考稳定性，并在双通道场景下辅助抑制公共模式干扰。虽然本研究的重点并不在于高密度医学级脑电前端设计，但对于双通道、低功耗、便携式采集设备而言，这种围绕输入稳定性和导联状态监测展开的前端设计，已经能够较好满足后续运动想象与基础生理特征观察的硬件需求。

\subsection{主控与无线传输电路设计}

主控与无线传输部分采用 nRF52832 作为核心芯片，负责完成双通道采样控制、前端参数配置、BLE 通信以及部分系统状态管理。图\ref{fig:c3_main_control}给出了主控局部原理图。该芯片集成 ARM Cortex-M4 内核、片上 SAADC 与 BLE 无线通信能力，适合构建低功耗、小体积的近体式采集终端。与“前端采样芯片 + 外部 MCU + 外部蓝牙模块”的分立式方案相比，采用单芯片主控可显著简化板级互连关系，降低系统面积与功耗开销。

从接口分配上看，系统将 `P0.02` 和 `P0.04` 配置为脑电采样输入通道，用于接收来自前端的两路模拟输出；`P0.18`、`P0.11`、`P0.17` 与 `P0.15` 分别构成与前端芯片的 SPI 控制链路；`P0.19`、`P0.20` 与 `P0.23` 则负责前端复位、使能与导联状态检测。此外，电路中配置了 `32\,\text{MHz}` 高频晶振和 `32.768\,\text{kHz}` 低频晶振，分别服务于主系统运行与低功耗时基需求；右侧射频匹配网络与天线连接则构成了 BLE 无线传输的硬件基础。

该部分电路设计的关键在于让同一主控芯片同时承担“采样入口”和“无线出口”两个角色。前者要求其能够以固定采样周期读取双通道 ADC 数据并及时写入缓存，后者要求其能够在蓝牙链路中完成 MTU 协商、通知发送和状态回执。nRF52832 的选型使这两项要求能够在单板平台上统一实现，从而为后续的软件架构设计提供了稳定硬件基础。

\begin{figure}[htbp]
    \centering
    \includegraphics[width=0.92\linewidth]{C3-2_Main_Control_Schematic.png}
    \caption{主控与无线传输局部原理图}
    \label{fig:c3_main_control}
\end{figure}

值得注意的是，主控电路中还预留了 I\textsuperscript{2}C 接口和中断引脚，用于接入基础生理信号扩展模块。结合现有固件和上位机实现，这一扩展链路可用于接入心率/血氧传感模块，从而在脑电采集之外增加辅助生理观测能力。考虑到本章重点在脑电硬件主链路设计，因此这里将其视为扩展能力而非本系统论述的核心。

\subsection{供电与接口设计}

为了满足便携式使用需求，系统在供电与接口部分采用了电池供电、Type-C 充电与板级稳压相结合的方案。根据器件清单与板级布局信息，系统包含 Type-C 接口、充电管理芯片、稳压器件、电池接口以及与按键相关的电源保持逻辑。该设计允许设备脱离外部实验台式电源独立运行，并为后续移动场景下的脑电采集提供硬件条件。

供电链路的设计重点并不在于追求高功率输出，而在于在较小板级面积内同时保证模拟前端、主控与无线链路的稳定供电。对于脑电采集系统而言，供电纹波、地参考波动以及蓝牙发射瞬态都有可能影响前端采集效果，因此在单板实现中将供电管理、电池接口和核心主控统一纳入设计，是实现整机稳定性的基础。此外，固件中还实现了基于按键与电源保持引脚的软关机逻辑，使设备具备更接近实际便携终端的工作方式，而非仅作为实验室内部的开发板使用。

\section{固件程序开发与流程设计}

在硬件电路基础上，系统固件负责将前端采集、主控处理、BLE 通信与状态管理整合为可运行的实时采集平台。本文固件基于 Zephyr RTOS 实现，采用模块化组织方式，将脑电采样、前端控制、数字滤波、无线传输以及扩展传感器处理等功能拆分到独立模块中。这样做的目的在于降低各功能之间的耦合度，并提高系统在后续功能增补和联调阶段的可维护性。

\subsection{固件核心工作流程}

系统上电后，固件首先完成主控基础外设初始化，包括 GPIO、ADC、蓝牙协议栈、NUS 服务以及脑电前端控制接口的配置；随后完成脑电前端芯片初始化与默认增益设置，并启动双通道采样模块。在采样任务正式运行后，系统以 `500\,\text{Hz}` 采样率驱动双通道采样流程，即每 `2\,\text{ms}` 触发一次采样调度。为了避免在定时器回调中执行耗时 ADC 读取操作，固件将真正的采样任务提交到独立工作队列中执行，从而降低实时调度被阻塞的风险。

在数据处理层面，采集到的双通道样本首先写入环形缓冲区，并附加设备侧时间戳；主循环随后从缓冲区中读取样本，执行数字滤波和状态处理，再根据当前 MTU 与批量打包策略构造 NUS 数据帧并通过 BLE 通知发送到上位机。同时，系统还持续监听来自上位机的文本命令，用于完成增益调整、滤波配置、缓存清理、状态查询和前端复位等控制任务。由此，系统形成了“周期采样 - 缓存入队 - 滤波处理 - 分帧发送 - 上位机回控”的闭环执行链路。

除脑电主链路外，固件中还包含独立的基础生理信号采集线程，用于读取扩展传感模块数据并通过文本消息向上位机回传结果。尽管这部分功能并非本章核心，但其存在说明该平台具备在低功耗嵌入式环境中同时管理多路采集任务与无线回传任务的能力。

\subsection{数字滤波与前端参数配置}

由于便携式脑电采集场景易受到工频干扰、低频漂移和高频噪声影响，本文在固件中实现了可配置的数字滤波策略。当前系统支持关闭滤波、工频陷波、低通滤波、组合滤波和带通滤波等多种模式，其中双通道默认采用组合滤波配置。相较于仅使用固定陷波器的单一实现方式，这种设计使系统能够根据不同实验任务和信号质量情况灵活切换处理方式。

从系统实现角度看，数字滤波的意义不仅在于改善波形观测效果，还在于将采集系统从“固定参数的示波终端”提升为“可配置的实验平台”。一方面，上位机可通过命令接口动态调整各通道滤波模式，使联调过程无需频繁修改固件；另一方面，前端增益也可通过命令接口在多个预设档位之间切换，从而适应不同被试个体差异和电极接触条件。脑电前端的导联脱落检测信号同样被纳入固件状态管理逻辑，并通过文本回执实时通知上位机。这种“前端状态感知 + 软件侧参数调节”的实现方式，有助于提高硬件系统在实际实验中的可用性。

\subsection{数据缓冲区与传输逻辑}

在无线双通道脑电采集中，采样任务与传输任务之间天然存在速率不一致问题。为解决这一问题，系统在固件中采用环形缓冲区作为采样与发送之间的解耦机制。当前实现中，缓冲区默认深度为 `1024` 个样本，最大支持 `2048` 个样本，并使用原子计数器记录已就绪样本数、总采样数以及缓冲区溢出次数。这样既能为上位机短时处理波动提供时间裕量，也便于后续对缓存状态和链路稳定性进行查询。

在传输侧，系统结合当前 BLE MTU 动态调整每帧携带的样本数量。连接建立后，设备会根据 MTU 更新结果重新计算批量发送的样本数上限，并为每个数据帧分配序列号。该策略相比逐点发送方式具有更高的有效负载利用率，更适合连续脑电流式数据的无线发送。与此同时，文本命令回执与二进制脑电帧采用同一 NUS 链路承载，但通过帧头魔数实现了解析层区分，从而在保持协议简洁的同时兼顾了实时数据流和状态控制流。总体而言，这种“环形缓冲 + 动态分帧 + 序列号跟踪”的实现机制，为双通道脑电在低功耗蓝牙链路中的稳定传输提供了软件支撑。

\section{系统联调验证与实现分析}

本章的重点在于系统设计与实现，因此本节不将上位机的实验效果和大规模系统性能评估作为主要内容，而是围绕目前已经完成的联调结果，对硬件系统的闭环可用性和实现特征进行分析。当前阶段已经能够确认的是：系统能够完成双通道脑电采样、前端参数配置、BLE 数据回传、导联状态上报以及上位机实时显示，说明该平台已经从单纯的电路设计方案推进到了可运行的系统原型阶段。

\subsection{脑电采集与显示功能验证}

从当前联调结果看，硬件系统已经能够实现双通道脑电信号的连续接收与实时显示。上位机界面如图\ref{fig:c3_ui}所示，能够同步显示双通道原始波形、频带能量分布、设备连接状态及部分扩展生理参数。这表明设备侧脑电采样、蓝牙传输与上位机解析链路已经完成闭环联调。尤其是导联状态和后端参数能够在界面中被直接观察与调整，说明系统不仅具备“采到数据”的能力，还具备“在采集过程中判断状态并完成配置”的能力。

对于第 3 章而言，这类联调验证的意义在于说明本系统不是停留在原理图和 PCB 层面的设计稿，而是已经实现了实际可运行的采集终端。虽然当前阶段尚未对特定实验范式下的分类表现进行展开讨论，但仅从波形连续显示、参数可调与状态可回读这三点来看，系统已经能够满足后续算法实验与基础生理特征观察的硬件支撑需求。

\subsection{链路时序与实时性分析}

在采样与显示的实时性方面，系统采用了从硬件到软件的分层时序控制策略。设备端以 `500\,\text{Hz}` 采样率执行双通道采样，对应的理论采样周期为 `2\,\text{ms}`；采样读取动作由工作队列完成，以避免阻塞定时器回调；传输侧则根据 MTU 动态调整批量打包样本数量，以减少蓝牙通知开销。上位机进一步采用实时波形刷新与较慢频谱更新相结合的策略，从而在保持主波形流畅显示的同时降低频域分析对主线程的压力。

需要说明的是，本阶段对实时性的分析主要基于链路实现方式和联调表现，而非基于示波器或统一测试工装得到的严格端到端延迟测量值。因此，本节更适合作为对系统实时性实现路径的分析，而不是对绝对时延指标的最终报告。即便如此，从当前的实现结构来看，采样、处理、传输与显示四个环节之间已经形成较为清晰的时序分工，这为后续进一步开展定量延迟测试奠定了基础。

\subsection{无线传输稳定性分析}

无线传输稳定性主要通过三个层面得到保障。首先，在设备端，脑电采样与发送之间通过环形缓冲区解耦，并对缓冲区使用量和溢出次数进行统计，从而为判断传输是否跟上采样速率提供依据；其次，在协议层，脑电帧引入了序列号与样本数量字段，上位机可据此识别潜在丢包与帧异常；最后，在控制层，系统支持状态查询、缓存清理、导联状态回报和错误回执，使联调过程中出现链路异常时能够快速定位问题。

从工程实现角度看，这种稳定性设计并不依赖单一的链路优化技巧，而是通过“采样缓冲 - 协议可观测 - 状态可查询”三个层面共同完成。对于便携式双通道脑电系统而言，这种设计比单纯追求链路峰值吞吐更有意义，因为其目标是支撑连续采集和可调试的实验平台，而不是极限数据传输系统。

\subsection{供电与便携化实现分析}

供电与便携化设计是该系统区别于单纯桌面开发板方案的重要特征。从板级器件和结构布局看，系统已经具备电池供电、Type-C 充电以及按键式上电/关机控制能力，说明其设计目标并非仅用于短时桌面调试，而是面向可移动、可穿戴或近体式实验场景。将供电管理纳入系统整体设计后，脑电采集、无线传输和参数控制能够在脱离外部有线供电条件下协同工作，这对于后续边缘化实验部署具有实际意义。

与此同时，当前阶段关于整机电流、续航时间和不同工作模式下功耗分布的量化测试仍有待进一步补充。因此，本节更多从实现角度说明系统已经具备便携终端的结构条件和供电逻辑，而不是给出最终的功耗评估结论。随着后续测试数据进一步完善，这部分内容还可以继续补充为更完整的系统性能分析。

\section{本章小结}

本章围绕双通道脑电采集硬件平台的设计与实现展开论述，完成了从总体架构、关键电路、固件流程到联调验证的系统性分析。本文所设计的平台以 KS1092 脑电前端和 nRF52832 主控为核心，构建了双通道脑电采集、可配置参数控制、BLE 无线传输和上位机实时显示的完整闭环；同时通过供电与接口设计使系统具备一定的便携部署能力。当前联调结果表明，该平台已经能够支撑双通道脑电信号的连续采样与实时可视化，为后续章节中的低通道算法验证和系统实验提供了硬件基础。
